\section*{Hoofdstuk \ref{ch:g11:reaction-energy} --- De energie van reacties: verbranding en bindingsenergieën}

\begin{solution}{exo:g11:reaction-energy:1}
Hout dat brandt: \omterm{def:g11:reaction-energy:exothermic}{exotherm}. IJs dat smelt: \omterm{def:g11:reaction-energy:exothermic}{endotherm} (het neemt warmte op
uit zijn omgeving, al is het geen \omterm{def:g7:chemical-reactions:reaction}{chemische reactie}). Koelkompres:
\omterm{def:g11:reaction-energy:exothermic}{endotherm}. Handwarmer: \omterm{def:g11:reaction-energy:exothermic}{exotherm}. Suiker maken in het zonlicht: \omterm{def:g11:reaction-energy:exothermic}{endotherm}
(de energie komt van het licht).
\end{solution}

\begin{solution}{exo:g11:reaction-energy:2}
$Q = 3.0 \times 890.7 = \qty{2.7e3}{kJ}$, ongeveer \qty{2.7}{MJ}.
\end{solution}

\begin{solution}{exo:g11:reaction-energy:3}
De \omterm{def:g7:chemical-reactions:reactant}{beginstoffen}. Het verschil wordt afgegeven aan de omgeving, vooral als
warmte.
\end{solution}

\begin{solution}{exo:g11:reaction-energy:4}
$n = 100 / 46.0 = \qty{2.17}{mol}$; $Q = 2.17 \times 1367.6 =
\qty{2.97e3}{kJ}$, ongeveer \qty{3.0}{MJ}.
\end{solution}

\begin{solution}{exo:g11:reaction-energy:5}
De energie die nodig is om één \omterm{def:g10:the-mole:amount}{mol} van zulke bindingen te verbreken,
waarbij \omterm{def:g7:atoms-and-molecules:molecule}{moleculen} en \omterm{def:g7:atoms-and-molecules:atom}{atomen} gasvormig zijn. Het gedeelde elektronenpaar
houdt de twee \omterm{def:g7:atoms-and-molecules:atom}{atomen} bij elkaar: ze uit elkaar trekken gaat tegen die
aantrekking in, en dat kost energie.
\end{solution}

\begin{solution}{exo:g11:reaction-energy:6}
Verbroken: $2 \times 436 + 498 = \qty{1370}{kJ}$; gevormd:
$4 \times 464 = \qty{1856}{kJ}$; $E_r \approx 1370 - 1856 =
\qty{-486}{kJ}$: \omterm{def:g11:reaction-energy:exothermic}{exotherm}.
\end{solution}

\begin{solution}{exo:g11:reaction-energy:7}
Ethanol \ce{CH3-CH2-O-H} heeft 5 \ce{C-H}, 1 \ce{C-C}, 1 \ce{C-O} en 1
\ce{O-H}. Verbroken: $5 \times 415 + 345 + 350 + 464 + 3 \times 498 =
\qty{4728}{kJ}$. Gevormd: $4 \times 741 + 6 \times 464 = \qty{5748}{kJ}$.
$E_r \approx \qty{-1020}{kJ}$, in grootte ongeveer een kwart kleiner dan
de gemeten \qty{1367.6}{kJ}.
\end{solution}

\begin{solution}{exo:g11:reaction-energy:8}
Propaan: $2219.2 / 0.0440 = \qty{5.04e4}{kJ/kg} = \qty{50.4}{MJ/kg}$.
Methaan: $890.7 / 0.0160 = \qty{55.7}{MJ/kg}$. Allebei zoals in het
diagram.
\end{solution}

\begin{solution}{exo:g11:reaction-energy:9}
$Q = 1500 \times 4.18 \times 85 = \qty{5.33e5}{J} = \qty{533}{kJ}$;
$n = 533 / 890.7 = \qty{0.598}{mol}$, dat is $0.598 \times 16.0 =
\qty{9.6}{g}$ methaan; als de helft van de energie verloren gaat,
\qty{19}{g}.
\end{solution}

\begin{solution}{exo:g11:reaction-energy:10}
Waterstof > methaan > propaan > octaan > ethanol.
$48.0 / 142.9 = \qty{0.34}{kg}$ waterstof.
\end{solution}

\begin{solution}{exo:g11:reaction-energy:11}
Verbroken: $946 + 3 \times 436 = \qty{2254}{kJ}$; gevormd:
$6 \times 390 = \qty{2340}{kJ}$; $E_r \approx \qty{-86}{kJ}$: licht
\omterm{def:g11:reaction-energy:exothermic}{exotherm}.
\end{solution}

\begin{solution}{exo:g11:reaction-energy:12}
Water: $250 \times 4.18 \times 20 = \qty{2.09e4}{J} = \qty{20.9}{kJ}$.
Ethanol: $1.20 / 46.0 = \qty{0.0261}{mol}$, die
$0.0261 \times 1367.6 = \qty{35.7}{kJ}$ zou kunnen afgeven. Rendement
$20.9 / 35.7 \approx \qty{59}{\percent}$. Verliezen: warmte die door de
\omterm{def:g4:air-a-mixture-of-gases:air}{lucht} wordt afgevoerd, warmte die het blikje en de driepoot opwarmt,
\omterm{def:g8:combustion-and-fuels:complete}{onvolledige verbranding} (roet), ethanol die onverbrand verdampt.
\end{solution}

\begin{solution}{exo:g11:reaction-energy:13}
\ce{2C8H18 + 25O2 -> 16CO2 + 18H2O}: 8 \omterm{def:g10:the-mole:amount}{mol} koolstofdioxide,
\qty{352}{g}, per \qty{5470}{kJ}, dus $352 / 5.470 = \qty{64.4}{g}$ per
megajoule. \ce{C3H8 + 5O2 -> 3CO2 + 4H2O}: \qty{132}{g} per
\qty{2219.2}{kJ}, dus \qty{59.5}{g} per megajoule. Methaan stoot met
\qty{49.4}{g} het minst uit.
\end{solution}

\begin{solution}{exo:g11:reaction-energy:14}
De tabel geeft gemiddelde bindingsenergieën, terwijl de
\ce{C=O}-bindingen van koolstofdioxide sterker zijn dan gemiddeld; de
gemeten waarden gelden voor vloeibaar water, waarvan de condensatie meer
energie afgeeft dan de reactie tussen gassen; en de methode behandelt
elke binding alsof ze los staat van haar buren. De schatting is toch
nuttig: ze geeft het juiste teken (\omterm{def:g11:reaction-energy:exothermic}{exotherm}) en de juiste orde van
grootte, zonder enige meting.
\end{solution}

\begin{solution}{exo:g11:reaction-energy:15}
Verbroken: $4 \times 464 = \qty{1856}{kJ}$; gevormd: $2 \times 436 + 498 =
\qty{1370}{kJ}$; $E_r \approx \qty{+486}{kJ}$: \omterm{def:g11:reaction-energy:exothermic}{endotherm}, dus er moet
energie worden toegevoerd. Als de waterstof later verbrandt, komt die
energie weer vrij: waterstof slaat energie op die elders is opgewekt, het
maakt er zelf geen.
\end{solution}

\begin{solution}{pb:g11:reaction-energy:1}
\textbf{1.} \ce{CH4 + 2O2 -> CO2 + 2H2O}.

\textbf{2.} \ce{C2H6O + 3O2 -> 2CO2 + 3H2O}.

\textbf{3.} \ce{2C8H18 + 25O2 -> 16CO2 + 18H2O}.

\textbf{4.} \ce{2H2 + O2 -> 2H2O}.

\textbf{5.} Waterstof.

\textbf{6.} \qty{16.0}{g/mol}, \qty{46.0}{g/mol}, \qty{114.0}{g/mol},
\qty{2.0}{g/mol}.

\textbf{7.} In \si{kJ} per \si{kg}, daarna \si{MJ/kg}: methaan
$890.7 / 0.0160$, \qty{55.7}{MJ/kg}; ethanol $1367.6 / 0.0460$,
\qty{29.7}{MJ/kg}; octaan $5470 / 0.1140$, \qty{48.0}{MJ/kg}; waterstof
$285.8 / 0.0020$, \qty{142.9}{MJ/kg}.

\textbf{8.} Waterstof > methaan > octaan > ethanol.

\textbf{9.} Waterstof is een heel licht gas: een kilogram ervan neemt bij
normale druk een enorm volume in, dus het moet in zware hogedruktanks
worden samengeperst.

\textbf{10.} Methaan: 4 \ce{C-H}; zuurstof: 2 \ce{O=O} verbroken.
Koolstofdioxide: 2 \ce{C=O}; water: $2 \times 2 = 4$ \ce{O-H} gevormd.

\textbf{11.} $4 \times 415 + 2 \times 498 = \qty{2656}{kJ}$.

\textbf{12.} $2 \times 741 + 4 \times 464 = \qty{3338}{kJ}$.

\textbf{13.} $E_r \approx 2656 - 3338 = \qty{-682}{kJ}$, tegen
\qty{-890.7}{kJ} gemeten: in grootte ongeveer \qty{23}{\percent} te
klein.

\textbf{14.} Gemiddelde bindingsenergieën (de \ce{C=O} van
koolstofdioxide is sterker dan gemiddeld); vloeibaar water bij de
meting, gassen bij de schatting.

\textbf{15.} $1000 / 890.7 = \qty{1.123}{mol}$.

\textbf{16.} $1.123 \times 44.0 = \qty{49.4}{g}$.

\textbf{17.} Ethanol: $1000 / 1367.6 = \qty{0.731}{mol}$, wat
\qty{1.462}{mol} koolstofdioxide geeft, \qty{64.3}{g}. Octaan:
$1000 / 5470 = \qty{0.183}{mol}$, wat \qty{1.463}{mol} geeft,
\qty{64.4}{g}.

\textbf{18.} Methaan: het heeft de meeste waterstofatomen per
koolstofatoom (4 tegen 2.25 bij octaan), en de \omterm{def:g4:what-a-fire-needs:burning}{verbranding} van waterstof
levert energie zonder koolstofdioxide.

\textbf{19.} Van de manier waarop de waterstof gemaakt wordt: uit water
met elektriciteit uit hernieuwbare bronnen stoot hij bijna geen
koolstofdioxide uit; gemaakt uit methaan stoot hij koolstofdioxide uit
in de fabriek in plaats van uit de uitlaat.

\textbf{20.} Methaan stoot ongeveer \qty{49}{g} koolstofdioxide per
megajoule uit.
\end{solution}
